\section{One charity, three questions}

The three readings agree that Christ's lordship gives love its object and measure, that grace enables obedience, and that the believer's good belongs within a people called to eternal life. Their differences concern where that common truth becomes most searching. The healing of the heart asks how a person who knows the command can become able to live it. The gathering of a people asks why those who receive the same gifts must bear one another across real differences. Faithful offering asks what becomes of prayer and reception when the worshipper returns to ordinary responsibilities.

The same proper consequently carries a different weight in each. The Introit exposes dependence on mercy, places a singular servant among blessed walkers, and begins the abandonment of self-defense. The Gospel heals the separation of love from obedience, places the body under its one Lord, and addresses the giver whose offering cannot reserve the rest of life. The Communion makes healing accountable, gathers humble gift-bearers around a common center, and demands that commitments be fulfilled. The Postcommunion joins all three: vices are cured in persons, the prayer is made together, and the desired fruit lasts beyond the action just celebrated.

The witnesses also differ in ways that enrich those questions. Bellarmine's created heavens and Augustine's spiritual heavens give different levels of meaning to the Gradual. Chrysostom's Spirit-given union and Thomas's ordered city emphasize vitality and structure. Jerome's alternatives concerning Daniel's confession preserve the difference between personal fault and representative solidarity. Bellarmine's Lord over mortal kings and Augustine's Lord over pride prevent judgment from becoming exclusively public or exclusively inward. These are not rival accounts of three different Masses. They show how the one formulary reaches the heart, the common body, and the life offered through worship.

The last word belongs to the petition for an eternal remedy. Love's present work is neither dismissed as impossible nor declared already finished. A people that asks to be healed can practice patient charity without pretending that its members have no wounds. A person who confesses the Lord can offer a faithful life without claiming that fidelity is self-created. The Mass commands, gives, judges, and heals within one approach to God.
